% ============================================================================
% Chapter 7 solutions: Simplex Method (dictionary form)
% Book source: part1-linear-programming/ch06-simplex/simplex-basis-driven.tex
% Exercise numbering follows \begin{ex} order in that file (17 exercises).
% All pivots and optima verified with exact-fraction basis algebra and
% scipy.optimize.linprog, July 2026.
% ============================================================================
\section{Chapter 7: Simplex Method}

This chapter's seventeen exercises track the chapter's arc: six warm-ups on
standard form and the initial dictionary (7.1 to 7.6), seven core simplex
solves in dictionary form including two Big-M problems (7.7 to 7.13), three
conceptual problems on optimality conditions, geometry of basic feasible
solutions, and the ratio test (7.14 to 7.16), and one challenge problem on
degeneracy from ratio-test ties (7.17). Every pivot below was verified by
exact-fraction basis computations and each optimum was confirmed with an LP
solver. Solutions printed in the book's Selected Solutions (7.5, 7.6, 7.7,
7.8, 7.9, 7.14, 7.17) are reproduced here, lightly edited, so that this
manual is self-contained.

\exsol{7.1}{Standard form and the first dictionary}{1}

\textbf{(a)} Add slack variables $s_1, s_2 \geq 0$:
\[
\begin{aligned}
\max \quad & 4x + 5y \\
\st \quad & x + 2y + s_1 = 12, \\
& 3x + 2y + s_2 = 24, \\
& x, y, s_1, s_2 \geq 0.
\end{aligned}
\]

\textbf{(b)} Solving each equation for its slack gives the initial dictionary
at the slack basis $\{s_1, s_2\}$:
\[
\begin{aligned}
\max \quad & z = 0 + 4x + 5y \\
\st \quad & s_1 = 12 - x - 2y, \\
& s_2 = 24 - 3x - 2y, \\
& x, y, s_1, s_2 \geq 0.
\end{aligned}
\]

\textbf{(c)} Setting the nonbasic variables $x = y = 0$ gives the basic
feasible solution $(x, y, s_1, s_2) = (0, 0, 12, 24)$ with objective value
$z = 0$.

\textbf{(d)} Under steepest ascent we pick the largest objective
coefficient: $5 > 4$, so $y$ enters first.

\exsol{7.2}{Convert to Standard Form}{1}

The objective is already a maximization and both variables are already
nonnegative. Add a slack variable $s_1 \geq 0$ to the $\leq$ constraint and
subtract a surplus variable $s_2 \geq 0$ from the $\geq$ constraint:
\[
\begin{aligned}
\max \quad & 3x_1 - 5x_2 \\
\st \quad & 2x_1 + x_2 + s_1 = 6, \\
& -x_1 + 3x_2 - s_2 = 2, \\
& x_1, x_2, s_1, s_2 \geq 0.
\end{aligned}
\]

\exsol{7.3}{Convert to Standard Form}{1}

Negate the objective to convert min to max: $\max\; x_1 - 2x_2 - x_3$.
Since $x_3$ is unrestricted, write $x_3 = x_3^+ - x_3^-$ with
$x_3^+, x_3^- \geq 0$. The first constraint is already an equality; add a
slack $s_1 \geq 0$ to the second:
\[
\begin{aligned}
\max \quad & x_1 - 2x_2 - x_3^+ + x_3^- \\
\st \quad & x_1 - x_2 + 2x_3^+ - 2x_3^- = 4, \\
& -2x_1 + x_2 + s_1 = 1, \\
& x_1, x_2, x_3^+, x_3^-, s_1 \geq 0.
\end{aligned}
\]

\exsol{7.4}{Convert to Standard Form}{1}

The objective is already a maximization. Handle the variables: $x_1$ is
unrestricted, so set $x_1 = x_1^+ - x_1^-$ with $x_1^+, x_1^- \geq 0$;
$x_2, x_3 \leq 0$, so set $x_2 = -x_2'$ and $x_3 = -x_3'$ with
$x_2', x_3' \geq 0$. Substituting, the objective becomes
$5x_1^+ - 5x_1^- - x_2' + 3x_3'$. Subtract a surplus $s_1 \geq 0$ from the
$\geq$ constraint; the equality stays as is:
\[
\begin{aligned}
\max \quad & 5x_1^+ - 5x_1^- - x_2' + 3x_3' \\
\st \quad & x_1^+ - x_1^- - 2x_2' + x_3' - s_1 = 7, \\
& -x_2' + 4x_3' = 3, \\
& x_1^+, x_1^-, x_2', x_3', s_1 \geq 0.
\end{aligned}
\]

\exsol{7.5}{Convert to Standard Form}{1}

(Book Selected Solution.) Negate the objective:
$\max\; -2x_1 - 4x_2$. Subtract a surplus $s_1 \geq 0$ from the $\geq$
constraint and add a slack $s_2 \geq 0$ to the $\leq$ constraint. Since
$x_1$ is unrestricted, replace $x_1 = x_1' - x_1''$ with
$x_1', x_1'' \geq 0$:
\[
\begin{aligned}
\max \quad & -2(x_1' - x_1'') - 4x_2 \\
\st \quad & (x_1' - x_1'') + x_2 - s_1 = 3, \\
& 2(x_1' - x_1'') - x_2 + s_2 = 5, \\
& x_1', x_1'', x_2, s_1, s_2 \geq 0.
\end{aligned}
\]

\exsol{7.6}{Convert to Standard Form}{1}

(Book Selected Solution.) Negate the objective:
$\max\; -5y_1 + 3y_2 - y_3$. Add a slack $s_1 \geq 0$ to the $\leq$
constraint; the equality needs no new variable. Since $y_3$ is
unrestricted, replace $y_3 = y_3' - y_3''$ with $y_3', y_3'' \geq 0$:
\[
\begin{aligned}
\max \quad & -5y_1 + 3y_2 - (y_3' - y_3'') \\
\st \quad & 2y_1 - y_2 + 3(y_3' - y_3'') = 7, \\
& -y_1 + 4y_2 + s_1 = 5, \\
& y_1, y_2, y_3', y_3'', s_1 \geq 0.
\end{aligned}
\]

\exsol{7.7}{Pivoting in a different direction}{2}

(Book Selected Solution, verified.) Entering $x$ first costs an extra
pivot compared with the steepest-ascent path: three pivots in all.

\textbf{Pivot 1: $x$ enters.} Ratio test: $9/1 = 9$ ($s_1$),
$16/2 = 8$ ($s_2$), $14/1 = 14$ ($s_3$); the minimum is $8$, so $s_2$
leaves. Solving $s_2$'s row for $x$ and substituting:
\[
\begin{aligned}
\max\quad & z = 16 + 2y - s_2 \\
\st\quad & s_1 = 1 - \tfrac12 y + \tfrac12 s_2, \\
& x = 8 - \tfrac12 y - \tfrac12 s_2, \\
& s_3 = 6 - \tfrac32 y + \tfrac12 s_2.
\end{aligned}
\]

\textbf{Pivot 2: $y$ enters} (only positive coefficient). Ratios:
$1/\tfrac12 = 2$ ($s_1$), $8/\tfrac12 = 16$ ($x$), $6/\tfrac32 = 4$
($s_3$); so $s_1$ leaves:
\[
\begin{aligned}
\max\quad & z = 20 - 4s_1 + s_2 \\
\st\quad & y = 2 - 2s_1 + s_2, \\
& x = 7 + s_1 - s_2, \\
& s_3 = 3 + 3s_1 - s_2.
\end{aligned}
\]

\textbf{Pivot 3: $s_2$ enters} (coefficient $+1$). The $y$ row has $s_2$
with a positive coefficient, so it imposes no limit; ratios $7/1 = 7$
($x$) and $3/1 = 3$ ($s_3$), so $s_3$ leaves:
\[
\begin{aligned}
\max\quad & z = 23 - s_1 - s_3 \\
\st\quad & x = 4 - 2s_1 + s_3, \\
& s_2 = 3 + 3s_1 - s_3, \\
& y = 5 + s_1 - s_3.
\end{aligned}
\]
All objective coefficients are negative, so this dictionary is optimal:
$(x, y) = (4, 5)$ with $z^* = 23$, the same optimum the steepest-ascent
path reaches in two pivots.

\exsol{7.8}{Simplex Method Practice}{2}

(Book Selected Solution, verified.) Initial dictionary:
\[
\begin{aligned}
\max\quad & z = 5x_1 + 3x_2 \\
\st\quad & s_1 = 12 - x_1 - x_2, \\
& s_2 = 16 - 2x_1 - x_2.
\end{aligned}
\]

\textbf{Pivot 1: $x_1$ enters} (coefficient $5$). Ratios: $12/1 = 12$ and
$16/2 = 8$, so $s_2$ leaves:
\[
\begin{aligned}
\max\quad & z = 40 + \tfrac12 x_2 - \tfrac52 s_2 \\
\st\quad & x_1 = 8 - \tfrac12 x_2 - \tfrac12 s_2, \\
& s_1 = 4 - \tfrac12 x_2 + \tfrac12 s_2.
\end{aligned}
\]

\textbf{Pivot 2: $x_2$ enters.} Ratios: $8/\tfrac12 = 16$ and
$4/\tfrac12 = 8$, so $s_1$ leaves:
\[
\begin{aligned}
\max\quad & z = 44 - s_1 - 2s_2 \\
\st\quad & x_2 = 8 - 2s_1 + s_2, \\
& x_1 = 4 + s_1 - s_2.
\end{aligned}
\]
Optimal: $(x_1, x_2) = (4, 8)$ with $z^* = 44$ (both constraints tight).

\exsol{7.9}{Simplex Method Practice}{2}

(Book Selected Solution, verified.) Initial dictionary:
\[
\begin{aligned}
\max\quad & z = 5x_1 + 8x_2 \\
\st\quad & s_1 = 30 - x_1 - 2x_2, \\
& s_2 = 30 - 3x_1 - x_2.
\end{aligned}
\]

\textbf{Pivot 1: $x_2$ enters} ($8 > 5$). Ratios: $30/2 = 15$ and
$30/1 = 30$, so $s_1$ leaves:
\[
\begin{aligned}
\max\quad & z = 120 + x_1 - 4s_1 \\
\st\quad & x_2 = 15 - \tfrac12 x_1 - \tfrac12 s_1, \\
& s_2 = 15 - \tfrac52 x_1 + \tfrac12 s_1.
\end{aligned}
\]

\textbf{Pivot 2: $x_1$ enters.} Ratios: $15/\tfrac12 = 30$ and
$15/\tfrac52 = 6$, so $s_2$ leaves:
\[
\begin{aligned}
\max\quad & z = 126 - \tfrac{19}{5} s_1 - \tfrac25 s_2 \\
\st\quad & x_1 = 6 + \tfrac15 s_1 - \tfrac25 s_2, \\
& x_2 = 12 - \tfrac35 s_1 + \tfrac15 s_2.
\end{aligned}
\]
Optimal: $(x_1, x_2) = (6, 12)$ with $z^* = 126$.

\exsol{7.10}{Simplex Method Practice}{2}

Add slacks $s_1, s_2$:
\[
\begin{aligned}
\max\quad & z = 2x_1 + 3x_2 + x_3 \\
\st\quad & s_1 = 32 - 4x_1 - 2x_2 - 5x_3, \\
& s_2 = 28 - 2x_1 - 4x_2 - 3x_3.
\end{aligned}
\]

\textbf{Pivot 1: $x_2$ enters} (largest coefficient $3$). Ratios:
$32/2 = 16$ and $28/4 = 7$, so $s_2$ leaves. With
$x_2 = 7 - \tfrac12 x_1 - \tfrac34 x_3 - \tfrac14 s_2$:
\[
\begin{aligned}
\max\quad & z = 21 + \tfrac12 x_1 - \tfrac54 x_3 - \tfrac34 s_2 \\
\st\quad & s_1 = 18 - 3x_1 - \tfrac72 x_3 + \tfrac12 s_2, \\
& x_2 = 7 - \tfrac12 x_1 - \tfrac34 x_3 - \tfrac14 s_2.
\end{aligned}
\]

\textbf{Pivot 2: $x_1$ enters} (coefficient $\tfrac12$). Ratios:
$18/3 = 6$ and $7/\tfrac12 = 14$, so $s_1$ leaves:
\[
\begin{aligned}
\max\quad & z = 24 - \tfrac{11}{6} x_3 - \tfrac16 s_1 - \tfrac23 s_2 \\
\st\quad & x_1 = 6 - \tfrac76 x_3 - \tfrac13 s_1 + \tfrac16 s_2, \\
& x_2 = 4 - \tfrac16 x_3 + \tfrac16 s_1 - \tfrac13 s_2.
\end{aligned}
\]
All objective coefficients are negative, so this is optimal:
$(x_1, x_2, x_3) = (6, 4, 0)$ with $z^* = 24$; both constraints are tight
and $x_3$ never enters.

\exsol{7.11}{Simplex Method Practice}{2}

Add slacks $s_1, s_2, s_3$:
\[
\begin{aligned}
\max\quad & z = 6x_1 + 8x_2 + 5x_3 \\
\st\quad & s_1 = 1800 - 4x_1 - x_2 - x_3, \\
& s_2 = 2000 - 2x_1 - 2x_2 - x_3, \\
& s_3 = 3200 - 4x_1 - 2x_2 - x_3.
\end{aligned}
\]

\textbf{Pivot 1: $x_2$ enters} (coefficient $8$). Ratios:
$1800/1 = 1800$, $2000/2 = 1000$, $3200/2 = 1600$, so $s_2$ leaves. With
$x_2 = 1000 - x_1 - \tfrac12 x_3 - \tfrac12 s_2$:
\[
\begin{aligned}
\max\quad & z = 8000 - 2x_1 + x_3 - 4s_2 \\
\st\quad & s_1 = 800 - 3x_1 - \tfrac12 x_3 + \tfrac12 s_2, \\
& x_2 = 1000 - x_1 - \tfrac12 x_3 - \tfrac12 s_2, \\
& s_3 = 1200 - 2x_1 + s_2.
\end{aligned}
\]

\textbf{Pivot 2: $x_3$ enters} (coefficient $+1$). The $s_3$ row does not
contain $x_3$; ratios $800/\tfrac12 = 1600$ ($s_1$) and
$1000/\tfrac12 = 2000$ ($x_2$), so $s_1$ leaves:
\[
\begin{aligned}
\max\quad & z = 9600 - 8x_1 - 2s_1 - 3s_2 \\
\st\quad & x_3 = 1600 - 6x_1 - 2s_1 + s_2, \\
& x_2 = 200 + 2x_1 + s_1 - s_2, \\
& s_3 = 1200 - 2x_1 + s_2.
\end{aligned}
\]
Optimal: $(x_1, x_2, x_3) = (0, 200, 1600)$ with $z^* = 9600$; the first
two constraints are tight, the third has slack $s_3 = 1200$, and $x_1$
stays out of the optimal solution.

\exsol{7.12}{Simplex Method Practice (Big-M)}{2}

Convert the minimization to a maximization, $\max\ z = -4x_1 - 6x_2 - 7x_3$,
subtract surplus variables $e_1, e_2 \geq 0$, and, since the origin is
infeasible, add artificial variables $a_1, a_2 \geq 0$ penalized by $-M$
in the objective. Solving for $a_1, a_2$ gives the starting dictionary
\[
\begin{aligned}
\max\quad & z = -50M + (2M - 4)x_1 + (3M - 6)x_2 + (3M - 7)x_3 - M e_1 - M e_2 \\
\st\quad & a_1 = 20 - x_1 - x_2 - 2x_3 + e_1, \\
& a_2 = 30 - x_1 - 2x_2 - x_3 + e_2.
\end{aligned}
\]

\textbf{Pivot 1: $x_2$ enters} (largest coefficient $3M - 6$). Ratios:
$20/1 = 20$ ($a_1$) and $30/2 = 15$ ($a_2$), so $a_2$ leaves:
\[
\begin{aligned}
\max\quad & z = (-90 - 5M) + \left(\tfrac{M}{2} - 1\right)x_1
  + \left(\tfrac{3M}{2} - 4\right)x_3 - M e_1
  + \left(\tfrac{M}{2} - 3\right)e_2 + \left(3 - \tfrac{3M}{2}\right)a_2 \\
\st\quad & a_1 = 5 - \tfrac12 x_1 - \tfrac32 x_3 + e_1 - \tfrac12 e_2 + \tfrac12 a_2, \\
& x_2 = 15 - \tfrac12 x_1 - \tfrac12 x_3 + \tfrac12 e_2 - \tfrac12 a_2.
\end{aligned}
\]

\textbf{Pivot 2: $x_3$ enters} (coefficient $\tfrac{3M}{2} - 4$, the
largest). Ratios: $5/\tfrac32 = \tfrac{10}{3}$ ($a_1$) and
$15/\tfrac12 = 30$ ($x_2$), so $a_1$ leaves. Both artificials are now
nonbasic and may be dropped:
\[
\begin{aligned}
\max\quad & z = -\tfrac{310}{3} + \tfrac13 x_1 - \tfrac83 e_1 - \tfrac53 e_2 \\
\st\quad & x_3 = \tfrac{10}{3} - \tfrac13 x_1 + \tfrac23 e_1 - \tfrac13 e_2, \\
& x_2 = \tfrac{40}{3} - \tfrac13 x_1 - \tfrac13 e_1 + \tfrac23 e_2.
\end{aligned}
\]

\textbf{Pivot 3: $x_1$ enters} (coefficient $\tfrac13$). Ratios:
$\tfrac{10}{3}/\tfrac13 = 10$ ($x_3$) and $\tfrac{40}{3}/\tfrac13 = 40$
($x_2$), so $x_3$ leaves:
\[
\begin{aligned}
\max\quad & z = -100 - x_3 - 2e_1 - 2e_2 \\
\st\quad & x_1 = 10 - 3x_3 + 2e_1 - e_2, \\
& x_2 = 10 + x_3 - e_1 + e_2.
\end{aligned}
\]
Optimal: $(x_1, x_2, x_3) = (10, 10, 0)$ with $\max z = -100$, that is,
minimum cost $z^* = 100$; both $\geq$ constraints are tight.

\begin{qaflag}
Tedium note: with the book's default largest-coefficient entering rule this
Big-M solve takes three pivots and passes through fractions in thirds,
which is right at the hand-computation limit. A shorter route exists:
entering $x_1$ first (its coefficient $2M - 4$ is also positive) drives
$a_1$ out immediately and reaches the same optimum in only two pivots.
An instructor may want to point students to that path, or the book could
add a hint such as ``entering $x_1$ first shortens the computation.''
The exercise itself is correct as stated.
\end{qaflag}

\exsol{7.13}{Simplex Method Practice (Big-M)}{2}

Convert to $\max\ z = -40x_1 - 48x_2 - 30x_3$, subtract surpluses
$e_1, e_2 \geq 0$, and add penalized artificials $a_1, a_2 \geq 0$:
\[
\begin{aligned}
\max\quad & z = -55M + (3M - 40)x_1 + (5M - 48)x_2 + (3M - 30)x_3 - M e_1 - M e_2 \\
\st\quad & a_1 = 25 - 2x_1 - 2x_2 - x_3 + e_1, \\
& a_2 = 30 - x_1 - 3x_2 - 2x_3 + e_2.
\end{aligned}
\]

\textbf{Pivot 1: $x_2$ enters} (largest coefficient $5M - 48$). Ratios:
$25/2 = 12.5$ ($a_1$) and $30/3 = 10$ ($a_2$), so $a_2$ leaves:
\[
\begin{aligned}
\max\quad & z = (-480 - 5M) + \left(\tfrac{4M}{3} - 24\right)x_1
  + \left(2 - \tfrac{M}{3}\right)x_3 - M e_1
  + \left(\tfrac{2M}{3} - 16\right)e_2 + \left(16 - \tfrac{5M}{3}\right)a_2 \\
\st\quad & a_1 = 5 - \tfrac43 x_1 + \tfrac13 x_3 + e_1 - \tfrac23 e_2 + \tfrac23 a_2, \\
& x_2 = 10 - \tfrac13 x_1 - \tfrac23 x_3 + \tfrac13 e_2 - \tfrac13 a_2.
\end{aligned}
\]

\textbf{Pivot 2: $x_1$ enters} ($\tfrac{4M}{3} - 24$ beats
$\tfrac{2M}{3} - 16$). Ratios: $5/\tfrac43 = \tfrac{15}{4}$ ($a_1$) and
$10/\tfrac13 = 30$ ($x_2$), so $a_1$ leaves. Both artificials are now
nonbasic; dropping them,
\[
\begin{aligned}
\max\quad & z = -570 - 4x_3 - 18e_1 - 4e_2 \\
\st\quad & x_1 = \tfrac{15}{4} + \tfrac14 x_3 + \tfrac34 e_1 - \tfrac12 e_2, \\
& x_2 = \tfrac{35}{4} - \tfrac34 x_3 - \tfrac14 e_1 + \tfrac12 e_2.
\end{aligned}
\]
All coefficients are negative, so this is optimal after just two pivots:
$(x_1, x_2, x_3) = \left(\tfrac{15}{4}, \tfrac{35}{4}, 0\right)$ with
minimum cost $z^* = 40 \cdot \tfrac{15}{4} + 48 \cdot \tfrac{35}{4}
= 150 + 420 = 570$. Both $\geq$ constraints are tight.

\begin{qaflag}
Judgment call: the optimal solution is fractional
$\left(\tfrac{15}{4}, \tfrac{35}{4}\right)$ even though all the data are
integers. The solve is short (two pivots) and the optimal value $570$ is
clean, so this is acceptable, but if an integer answer is preferred,
changing the first right-hand side from $25$ to $24$ yields the integral
optimum $(3, 9, 0)$ with cost $552$ (verified).
\end{qaflag}

\exsol{7.14}{Dictionary Optimality}{2}

(Book Selected Solution, verified.) The basic solution reads off the
constant terms: $x_2 = 1$, $x_4 = 2$, $x_3 = b$, with
$x_1 = x_5 = x_6 = 0$.
\begin{enumerate}
\item \textbf{Feasible and optimal:} feasibility requires every constant
term nonnegative, so $b \geq 0$; optimality requires no improving
direction, so $c_1 \leq 0$, $c_5 \leq 0$, $c_6 \leq 0$. The values
$a_1, a_2$ may be arbitrary.
\item \textbf{Not feasible:} $b < 0$ (the constants $1$ and $2$ are
already nonnegative, so only the $x_3$ row can fail).
\item \textbf{Degenerate:} $b = 0$, so that the basic variable $x_3$ sits
at value zero (feasibility of the rest is automatic).
\item \textbf{Feasible and unbounded:} we need $b \geq 0$ together with an
entering variable whose objective coefficient is positive and whose
increase never decreases any basic variable. The variable $x_5$ appears
with fixed coefficient $-1$ in the $x_3$ row and $x_6$ with $-1$ in both
the $x_4$ and $x_3$ rows, so the ratio test always stops them. Only $x_1$
can certify unboundedness: its column coefficients are $+1$, $+1$, and
$-a_2$, so the conditions are $c_1 > 0$ and $a_2 \leq 0$ (with $b \geq 0$
for feasibility).
\end{enumerate}

\exsol{7.15}{Basic Feasible Solutions}{2}

\begin{qaflag}
Book fix applied: in the source, all but one constraint of this exercise
had been commented out, and the one visible constraint
($-x_1 + x_2 \leq 2$) did not match the figure. The constraint system was
reconstructed exactly from the figure (all five labeled vertices check
out: $A = (2,1)$, $B = (5,1)$, $C = (7,3)$, $D = (0, \tfrac{16}{3})$,
$E = (0, \tfrac{7}{3})$) and the book statement was repaired to
\[
\max\ 2x_1 + x_2 \quad \st\quad 2x_1 + 3x_2 \geq 7,\ \
x_2 \geq 1,\ \ x_1 - x_2 \leq 4,\ \ x_1 + 3x_2 \leq 16,\ \ x_1, x_2 \geq 0,
\]
with surplus variables $s_1, s_2$ on the two $\geq$ constraints and slacks
$s_3, s_4$ on the two $\leq$ constraints. The solution below matches the
repaired statement.
\end{qaflag}

\textbf{(a)} The vertex $D = (0, \tfrac{16}{3})$ is the intersection of
the boundaries $x_1 = 0$ and $x_1 + 3x_2 = 16$ (that is, $s_4 = 0$). So
the nonbasic variables are $N = \{x_1, s_4\}$ and the basic variables are
$B = \{x_2, s_1, s_2, s_3\}$, with values
\[
x_2 = \tfrac{16}{3}, \qquad s_1 = 9, \qquad s_2 = \tfrac{13}{3},
\qquad s_3 = \tfrac{28}{3}.
\]

\textbf{(b)} Moving from $D$ to $C = (7, 3)$ travels along the edge
$s_4 = 0$ while $x_1$ increases from $0$: so \emph{$x_1$ enters}. At $C$
the boundary $x_1 - x_2 = 4$ becomes tight, i.e., $s_3 = 0$: so
\emph{$s_3$ leaves}. This agrees with the ratio test in the dictionary at
$D$, where increasing $x_1$ (with $s_4 = 0$) drives
$s_3 = \tfrac{28}{3} - \tfrac43 x_1$ to zero first, at $x_1 = 7$,
before $s_2$ (at $x_1 = 13$) or $x_2$ (at $x_1 = 16$).

\exsol{7.16}{What the ratio test protects}{2}

\textbf{(a)} Pivoting on the $s_2$ row means solving
$s_2 = 9 - \tfrac32 x + \tfrac12 s_3$ for $x$, giving
$x = 6 - \tfrac23 s_2 + \tfrac13 s_3$, and substituting everywhere:
\[
\begin{aligned}
\max\quad & z = 24 - \tfrac13 s_2 - \tfrac43 s_3 \\
\st\quad & s_1 = -1 + \tfrac13 s_2 + \tfrac13 s_3, \\
& x = 6 - \tfrac23 s_2 + \tfrac13 s_3, \\
& y = 4 + \tfrac13 s_2 - \tfrac23 s_3.
\end{aligned}
\]
The basic solution is $(x, y, s_1, s_2, s_3) = (6, 4, -1, 0, 0)$: the
basic variable $s_1$ is negative, so the dictionary is infeasible. (Indeed
the point $(x, y) = (6, 4)$ violates $x + y \leq 9$.) Note that the
``objective value'' $24$ exceeds the true optimum $23$ precisely because
the point is outside the feasible region.

\textbf{(b)} As the entering variable $x$ increases, every basic variable
with a negative $x$-coefficient decreases toward zero; the ratio of each
row measures how far $x$ can go before that row's basic variable hits
zero. The minimum ratio is the first blocking value. Pivoting on any
other row pushes $x$ past that first block, so the basic variable that
owned the minimum ratio (here $s_1$, blocked at $x = 4$) is driven
negative. The minimum-ratio rule protects exactly the nonnegativity of
all the basic variables that remain in the basis.

\exsol{7.17}{A tie in the ratio test forces degeneracy}{3}

(Book Selected Solution, verified, with part (c) restated.)

\textbf{(a)} The initial dictionary is
\[
\begin{aligned}
\max\quad & z = 2x + y \\
\st\quad & s_1 = 4 - x - y, \\
& s_2 = 8 - 2x - y.
\end{aligned}
\]
With $x$ entering (and $y = 0$), the ratio test gives $4/1 = 4$ for $s_1$
and $8/2 = 4$ for $s_2$: a tie at $x = 4$.

\textbf{(b)} Let $s_1$ leave. Then $x = 4 - y - s_1$ and
\[
\begin{aligned}
\max\quad & z = 8 - y - 2s_1 \\
\st\quad & x = 4 - y - s_1, \\
& s_2 = 0 + y + 2s_1.
\end{aligned}
\]
The basic variable $s_2$ has value $0$: the dictionary is degenerate. It
is also optimal (both reduced costs negative), with $(x, y) = (4, 0)$ and
$z^* = 8$.

\textbf{(c)} Let $s_2$ leave instead. Then
$x = 4 - \tfrac12 y - \tfrac12 s_2$ and
\[
\begin{aligned}
\max\quad & z = 8 + 0 \cdot y - s_2 \\
\st\quad & x = 4 - \tfrac12 y - \tfrac12 s_2, \\
& s_1 = 0 - \tfrac12 y + \tfrac12 s_2.
\end{aligned}
\]
Again a basic variable ($s_1$) equals zero: degenerate. The reduced cost
of $y$ is exactly $0$. A zero reduced cost normally signals alternative
optima along the corresponding edge; here the degenerate row
$s_1 = 0 - \tfrac12 y + \cdots$ blocks $y$ at once, so the optimal
\emph{solution} $(4, 0)$ is unique, but there are multiple optimal
\emph{bases} describing it. This is the typical fingerprint of degeneracy.

\textbf{(d)} If two rows tie at the minimum ratio, the entering variable
stops at a value where both corresponding basic variables hit zero
simultaneously. Only one of them leaves the basis; the other remains
basic at value zero, so the new dictionary necessarily has a basic
variable equal to zero, which is the definition of degeneracy. A
construction: $\max\ x + y$ subject to $x \leq 3$, $x + 2y \leq 3$,
$x, y \geq 0$; entering $x$ gives ratios $3/1 = 3$ and $3/1 = 3$, a tie.
